\uuid{RAFV}
\exo7id{3954}
\titre{exo7 3954}
\auteur{quercia}
\organisation{exo7}
\datecreate{2010-03-11}
\isIndication{false}
\isCorrection{false}
\chapitre{Dérivabilité des fonctions réelles}
\sousChapitre{Autre}
\cpge{mpsi-mp2i.fonctions-valeurs-reelles-complexes.b, pcsi.techniques-calcul-fonctions.b, ptsi.techniques-calcul-fonctions.b, mpsi-mp2i.polynomes-fractions-rationnelles.c, pcsi.polynomes.c, ptsi.polynomes.c, bcpst1.derivation.b, bcpst1.polynomes-reels.b}
\module{Analyse}
\niveau{L1}
\difficulte{}

\contenu{
\question{
Soit $n \in \N,\ n \ge 2$, et $a,b \in \R$.
Montrer que l'équation $x^n + ax + b = 0$ ne peut avoir plus de deux racines réelles
distinctes si $n$ est pair, et plus de trois racines réelles distinctes si
$n$ est impair.
}
}
